% TEMPLATE-MILC-v3.tex - plantilla maestra UNICA de las guias MILC v3.
%
% La usan las tres familias de guias, cada una con su driver:
%   · Grados 6-11  -> scripts/build-guias-g11.py
%   · Bebras       -> scripts/build-guias-bebras.py
%   · Semillero    -> scripts/build-guias-semillero.py
%
% Antes habia tres copias casi identicas (~400 lineas iguales, 6 distintas):
% cada mejora habia que aplicarla tres veces, y en la practica no se hacia
% (el motor de recursos vivia solo en dos de ellas). Lo que varia entre
% familias esta parametrizado en 6 placeholders que cada driver llena
% (se nombran aqui SIN su sintaxis real: el builder reemplaza por texto plano,
% tambien dentro de los comentarios, y un valor multilinea rompe el preambulo):
%
%   HEADER_IZQ        encabezado izquierdo de pagina
%   HEADER_DER        encabezado derecho de pagina
%   PORTADA_ETIQUETA  rotulo sobre el numero grande (GRADO/MOMENTO/MODULO)
%   PORTADA_SERIE     linea de serie de la portada
%   PORTADA_ID        identificador de la guia en la portada
%   PORTADA_CIRCULO   el circulito de grado (°), o vacio si no aplica
%
% El resto de claves las documenta content/guias/_SCHEMA.md. El script valida
% que no queden placeholders sin reemplazar antes de escribir el archivo final.

\documentclass[11pt,letterpaper]{article}
\usepackage[margin=1.75cm]{geometry}
\usepackage[table]{xcolor}
\usepackage{fontspec}
\usepackage[spanish]{babel}
\usepackage{tikz}
% Librerías TikZ para los diagramas de las guías (bloque `recursos:`):
%  · babel  → babel en español vuelve activos caracteres como «>», y TikZ los usa
%             en sus opciones (>=stealth, ->). Sin esta librería, cualquier
%             diagrama con flechas revienta con «\language@active@arg> has an
%             extra }». Debe ir ANTES de que se dibuje nada.
%  · positioning        → sintaxis «below left=2pt and 3pt».
%  · decorations.pathreplacing → llaves ({brace}) para señalar rangos.
\usetikzlibrary{babel,positioning,decorations.pathreplacing}
\usepackage{graphicx}
% Circuitos: el trabajo de electrónica se hace hoy con micro:bit y sus sensores
% integrados (MakeCode, sin cableado), así que no se necesitan esquemáticos.
% Si algún día entran montajes con protoboard, basta descomentar esta línea:
% los diagramas .tex/.tikz declarados en `recursos:` ya se incrustan solos.
% \usepackage{circuitikz}
\usepackage[most]{tcolorbox}
\usepackage{tabularx}
\usepackage{array}
\usepackage{fancyhdr}
\usepackage{titlesec}
\usepackage{ragged2e}
\usepackage{enumitem}
\usepackage{microtype}

% Helvetica Neue no trae la flecha U+2192, asi que cada «→» escrito literalmente
% en el YAML se compilaba a NADA, en silencio (462 flechas en 61 guias: rutas del
% tipo «paleta → tipografia → lockup» salian rotas). Mapeamos el caracter a su
% version matematica, que si tiene glifo. Asi el autor puede seguir escribiendo
% «→» comodamente en el YAML.
\usepackage{newunicodechar}
\newunicodechar{→}{\ensuremath{\rightarrow}}

\setmainfont{Helvetica Neue}
\setsansfont{Helvetica Neue}
\setlength{\parindent}{0pt}
\setlength{\parskip}{6pt}
\hyphenpenalty=200
\exhyphenpenalty=200
\pretolerance=400
\tolerance=2000
\setlength{\emergencystretch}{1em}
\renewcommand{\arraystretch}{1.25}
\setlist{leftmargin=5mm,itemsep=1mm,topsep=1mm}
\newcolumntype{Y}{>{\RaggedRight\arraybackslash}X}

\definecolor{milcMagenta}{HTML}{E6007E}
\definecolor{milcVino}{HTML}{5A0038}
\definecolor{milcTurquesa}{HTML}{0093A5}
\definecolor{milcVerde}{HTML}{58A923}
\definecolor{milcMostaza}{HTML}{E5B400}
\definecolor{milcOcre}{HTML}{B86600}
\definecolor{milcNegro}{HTML}{241816}
\definecolor{milcGris}{HTML}{F7F7F4}
\definecolor{milcLinea}{HTML}{E8E2DD}
\definecolor{milcGradePrimary}{HTML}{0066FF}
\definecolor{milcGradeDark}{HTML}{003D99}
\definecolor{milcGradeSoft}{HTML}{D6E8FF}
\definecolor{milcGradeAccent}{HTML}{CFE7FF}
\definecolor{milcGradeText}{HTML}{FFFFFF}
\color{milcNegro}

\pagestyle{fancy}
\fancyhf{}
\lhead{\small Tecnología e Informática · Grado 11}
\rhead{\small Guía 1 · MILC v3}
\cfoot{\small\thepage}
\renewcommand{\headrulewidth}{0pt}

\newtcolorbox{titlebox}[1]{
  enhanced,colback=#1,colframe=#1,arc=7mm,boxrule=0pt,
  left=7mm,right=7mm,top=3mm,bottom=3mm,
  before skip=8pt,after skip=8pt,
  fontupper=\sffamily\bfseries\Large\color{white}
}

\newtcolorbox{softbox}[3]{
  enhanced,colback=#2,colframe=#1,borderline west={4pt}{0pt}{#1},
  arc=2mm,boxrule=.4pt,left=4mm,right=4mm,top=3mm,bottom=3mm,
  before skip=6pt,after skip=8pt,title={#3},coltitle=white,
  colbacktitle=#1,fonttitle=\sffamily\bfseries,fontupper=\RaggedRight,
  attach boxed title to top left={xshift=3mm,yshift=-2mm},
  boxed title style={arc=2mm, boxrule=0pt}
}

\newtcolorbox{drawbox}[2]{
  enhanced,colback=white,colframe=#1,arc=4mm,boxrule=1pt,
  left=5mm,right=5mm,top=4mm,bottom=4mm,before skip=8pt,after skip=8pt,
  title={#2},coltitle=white,colbacktitle=#1,fonttitle=\sffamily\bfseries,
  fontupper=\RaggedRight,
  attach boxed title to top left={xshift=4mm,yshift=-2mm},
  boxed title style={arc=3mm, boxrule=0pt}
}

\newtcolorbox{infoband}[1]{
  enhanced,colback=#1!8,colframe=#1!50,arc=2mm,boxrule=.5pt,
  left=4mm,right=4mm,top=2mm,bottom=2mm,before skip=4pt,after skip=4pt,
  fontupper=\small\RaggedRight
}

% Puente narrativo entre secciones (contrato editorial Conéctate).
% Da continuidad: "Acabas de X. Ahora vas a Y porque Z."
% Visualmente: banda izquierda en color del grado, texto en itálica pequeña.
\newtcolorbox{puentebox}{
  enhanced,colback=milcGradeSoft,colframe=milcGradeDark,
  borderline west={3pt}{0pt}{milcGradeDark},
  arc=0mm,boxrule=0pt,left=5mm,right=4mm,top=2.5mm,bottom=2.5mm,
  before skip=4pt,after skip=8pt,
  fontupper=\itshape\small\color{milcNegro}\RaggedRight
}
% La flecha va en modo matematico a proposito: Helvetica Neue NO tiene el glifo
% U+2192, asi que \textrightarrow se compilaba a NADA («Missing character: There
% is no → in font Helvetica Neue») y el puente salia sin flecha. En modo
% matematico la flecha viene de la fuente matematica y si se dibuja.
\newcommand{\puente}[1]{%
  \begin{puentebox}%
    \textcolor{milcGradeDark}{$\rightarrow$}\hspace{2mm}#1%
  \end{puentebox}%
}

\newcommand{\linead}[1]{\noindent\color{milcLinea}\rule{#1}{0.5pt}\color{milcNegro}}
\newcommand{\respuesta}[1]{\par\vspace{2mm}\foreach \n in {1,...,#1}{\noindent\color{milcLinea}\rule{\linewidth}{0.5pt}\par\vspace{5mm}}\color{milcNegro}}
\newcommand{\checkbox}{\textcolor{milcGradeDark}{\fbox{\phantom{x}}}\hspace{5pt}}

\newcommand{\phasecard}[4]{
\begin{tcolorbox}[enhanced,colback=#3,colframe=#2,arc=3mm,boxrule=.5pt,height=3.05cm,valign=center,halign=center,left=1.5mm,right=1.5mm,top=2mm,bottom=2mm]
{\sffamily\bfseries\fontsize{22}{24}\selectfont\textcolor{#2}{#1}}\par
{\sffamily\bfseries\small #4}
\end{tcolorbox}
}

% ── Recursos visuales de la guía ──────────────────────────────────────
% Declarados en el bloque `recursos:` del YAML y emitidos por el builder.
% \guiaFigura   → imagen rasterizada o PDF (foto, carta celeste, montaje).
% \guiaDiagrama → fuente TikZ/circuitikz (.tex) incrustada como vector:
%                 es la vía para circuitos y esquemáticos editables.
% #1 = ruta relativa al .tex · #2 = pie (puede ir vacío)
\newcommand{\guiaFigura}[2]{%
\begin{center}
\includegraphics[width=0.88\linewidth,height=0.38\textheight,keepaspectratio]{#1}\\[1.5mm]
{\footnotesize\itshape\textcolor{milcNegro!75}{#2}}
\end{center}
}

\newcommand{\guiaDiagrama}[2]{%
\begin{center}
\input{#1}\\[1.5mm]
{\footnotesize\itshape\textcolor{milcNegro!75}{#2}}
\end{center}
}

\begin{document}
\thispagestyle{empty}
\begin{tikzpicture}[remember picture,overlay]
  \fill[white] (current page.south west) rectangle (current page.north east);
  \path[fill=milcGradePrimary,rounded corners=16mm]
    ([xshift=.55cm,yshift=-.55cm]current page.north west) rectangle
    ([xshift=-.55cm,yshift=3.65cm]current page.south east);

  \node[anchor=north west,text=milcGradeText] at ([xshift=2.0cm,yshift=-1.85cm]current page.north west)
    {\sffamily\bfseries\fontsize{14}{16}\selectfont GRADO};
  \node[anchor=north west,text=milcGradeText,opacity=.82] at ([xshift=2.0cm,yshift=-2.75cm]current page.north west)
    {\sffamily\bfseries\fontsize{9}{11}\selectfont SERIE GUÍAS · TECNOLOGÍA E INFORMÁTICA};

  \begin{scope}[shift={([xshift=-3.05cm,yshift=-2.55cm]current page.north east)}]
    \fill[white,opacity=.80] (-.62,-.52) rectangle (.62,.52);
    \draw[milcGradeText,line width=1pt] (-.62,-.52) rectangle (.62,.52);
    \fill[milcGradeDark] (-.42,-.38) rectangle (-.22,.06);
    \fill[milcGradeAccent] (-.10,-.38) rectangle (.10,.24);
    \fill[milcGradePrimary!60!black] (.22,-.38) rectangle (.42,.38);
  \end{scope}

  \node[anchor=west,text=milcGradeText] at ([xshift=1.9cm,yshift=-12.85cm]current page.north west)
    {\sffamily\bfseries\fontsize{124}{124}\selectfont 11};
  % El circulito de grado (°) solo aplica a las guias de grado («11°»); en
  % Bebras y Semillero el numero grande es un momento/modulo, no un grado.
  \draw[milcGradeText,line width=1.7pt]
    ([xshift=4.52cm,yshift=-11.68cm]current page.north west) circle (.13cm);

  \node[anchor=west,text=milcGradeText,text width=8.7cm,align=left] at ([xshift=10.35cm,yshift=-11.6cm]current page.north west)
    {
     {\sffamily\bfseries\fontsize{19}{22}\selectfont Mi marca digital ---\\del oficio heredado\\a la huella propia}\\[4mm]
     {\sffamily\bfseries\fontsize{23}{26}\selectfont Guía 1-11-TIC}\\[2mm]
     {\sffamily\fontsize{13}{16}\selectfont Período 1 · Presencia y marca digital}\\[5mm]
     {\sffamily\bfseries\fontsize{11}{13}\selectfont Institución Educativa Sor María Juliana}\\[1mm]
     {\sffamily\fontsize{10}{12}\selectfont Autor: Dr. Álvaro Cárdenas Orozco}
    };

  \path[fill=milcGradeSoft,rounded corners=7mm]
    ([xshift=1.35cm,yshift=.85cm]current page.south west) rectangle
    ([xshift=-1.35cm,yshift=4.25cm]current page.south east);
  \draw[milcGradeDark,line width=.65pt,rounded corners=7mm]
    ([xshift=1.35cm,yshift=.85cm]current page.south west) rectangle
    ([xshift=-1.35cm,yshift=4.25cm]current page.south east);
  \node[anchor=center,text=milcNegro,text width=17.0cm,align=center] at ([yshift=2.55cm]current page.south)
    {\sffamily\fontsize{10}{12}\selectfont Metodología MILC · Apertura ancestral · Pensamiento computacional · Triángulo Dussel-Estoicismo-Floridi\\
    \textcolor{milcGradeDark}{\bfseries Producto:} Manifiesto de marca personal en una página de cuaderno};
\end{tikzpicture}
\mbox{}
\newpage

\section*{Apertura: Mi marca digital --- del oficio heredado a la huella propia}

\begin{softbox}{milcGradeDark}{milcGradeSoft}{Saber ancestral}
En el Valle del Cauca, el Pacífico y el campo colombiano, una persona se conocía por dos cosas: su \textbf{oficio} y su \textbf{palabra}. ``Es el hijo del panadero'', ``los Cárdenas son maestros'', ``esa familia es seria''. Una marca nunca nació en internet --- nació en la plaza del pueblo, en la voz que pasa de vecino a vecino. La marca, en el fondo, es \emph{la promesa que tu nombre sostiene cuando no estás presente}.
\end{softbox}

\begin{softbox}{milcTurquesa}{milcGris}{Saber contemporáneo}
Hoy esa promesa se construye en redes, perfiles, búsquedas de Google y huella digital. Un futuro empleador, un jurado de universidad o un cliente potencial te lee \emph{antes} de conocerte. Tu marca personal trabaja por ti --- o en tu contra --- 24 horas al día. Construirla con intención hoy te ahorra borrar errores mañana.
\end{softbox}

\begin{softbox}{milcMostaza}{milcGris}{Pregunta-puente}
\textit{¿Qué prometía el nombre de tu familia en su barrio o vereda hace 50 años? ¿Qué promete hoy en redes el nombre que tú llevas?}
\end{softbox}

\begin{softbox}{milcVerde}{milcGris}{Saber y saber hacer}
Al terminar esta guía podrás: (1) \textbf{identificar} qué huella digital tienes hoy; (2) \textbf{explicar} tu propuesta de valor en una frase concreta; (3) \textbf{crear} tu manifiesto de marca personal con cinco bloques fijos; (4) \textbf{evaluar} con criterios el manifiesto de un compañero.
\end{softbox}



\small
\begin{tabularx}{\textwidth}{>{\bfseries}p{3.0cm}Y}
\rowcolor{milcGradePrimary!12}Autor & Dr. Álvaro Cárdenas Orozco\\
DBA & Comunicación profesional en entornos digitales (MEN, Lineamientos T\&I)\\
Referentes & Floridi (infoética) · Dussel (estética de la liberación) · Estoicismo\\
Producto final & Manifiesto de marca personal en una página de cuaderno\\
\end{tabularx}
\normalsize

\puente{Antes de empezar, mira el viaje completo. Cinco estaciones, una promesa: al final de la guía habrás escrito --- de tu propia mano --- el manifiesto de marca personal que va a ser tu ancla durante todo el periodo.}

\begin{titlebox}{milcGradeDark}
Ruta MILC: así vamos a aprender
\end{titlebox}
\begin{tabularx}{\textwidth}{>{\centering\arraybackslash}X>{\centering\arraybackslash}X>{\centering\arraybackslash}X>{\centering\arraybackslash}X}
\phasecard{1}{milcVerde}{milcGris}{Escucha\\[-1mm]\small Observo, escucho y pregunto.} &
\phasecard{2}{milcTurquesa}{milcGris}{Sistematización\\[-1mm]\small Pensamiento computacional.} &
\phasecard{3}{milcMagenta}{milcGris}{Praxis\\[-1mm]\small Construyo y aplico.} &
\phasecard{4}{milcVino}{milcGris}{Evaluación\\[-1mm]\small Reflexión liberadora.}
\end{tabularx}


\puente{Una marca no se construye al aire: se levanta sobre lo que ya existe. Antes de proyectar quién quieres ser, vamos a mirar la huella que ya dejaste en internet. Es el espejo más duro y el más útil que tienes hoy.}

\begin{titlebox}{milcVerde}
Fase 1 · Escucha
\end{titlebox}

\begin{softbox}{milcVerde}{milcGris}{Escena para escuchar}
\textbf{Actividad 1 · ANALIZA --- Audita tu huella digital actual} (20 min · individual). En tu celular, abre Google en una pestaña de incógnito y busca tu nombre completo entre comillas. Revisa los primeros 10 resultados de Google y los primeros 5 de Google Imágenes. Luego revisa las últimas 20 publicaciones de tus 3 redes principales (Instagram, TikTok, Facebook o las que uses) con ojo de futuro empleador.
\end{softbox}

\checkbox Busqué mi nombre completo en Google incógnito y revisé los primeros 10 resultados.\\
\checkbox Revisé las últimas 20 publicaciones de mis 3 redes principales con ojo de futuro empleador.\\
\checkbox Clasifiqué cada elemento en: suma, neutro o resta a mi marca.

\begin{infoband}{milcVerde}


\textbf{Lo que va al cuaderno (Actividad 1).} Título: «Actividad 1 --- Auditoría de mi huella digital». \textbf{Formato:} tabla de 3 columnas (Suma~|~Neutro~|~Resta), mínimo 3 elementos por columna, más un párrafo de diagnóstico de 5--8 renglones. \textbf{Sabes que terminaste cuando} cada elemento es específico (URL, post, foto, comentario --- no genérico) y el párrafo responde con honestidad, reconociendo al menos un sesgo de tu propia mirada.
\end{infoband}

\puente{Ya viste lo que la red dice de ti hoy. Ahora vamos a entender qué hace fuerte a una marca, para que la tuya tenga \textbf{criterio}, no \textbf{instinto}. Cuatro pilares del pensamiento computacional aplicados a una sola pregunta: ¿qué prometo, a quién, y por qué?}

\begin{titlebox}{milcTurquesa}
Fase 2 · Sistematización con pensamiento computacional
\end{titlebox}

\begin{softbox}{milcTurquesa}{milcGris}{Cuatro pilares aplicados al tema}
Una marca personal se compone de cuatro elementos. Vamos a descomponerlos con los pilares del pensamiento computacional para construir tu \textbf{propuesta de valor} con criterio --- no por instinto ni por copia.
\end{softbox}

\small
\begin{tabularx}{\textwidth}{>{\bfseries\RaggedRight\arraybackslash}p{3.6cm}Y}
\rowcolor{milcTurquesa!90}\color{white}Pilar & \color{white}\bfseries Aplicación al tema\\
1. Descomposición & \textbf{Descomposición.} Separamos la marca en sus 4 piezas para no diluirla: \textbf{identidad} (quién soy hoy, no quién quisiera ser), \textbf{audiencia} (a quién hablo, lo más concreto posible: ``estudiantes de mi universidad que buscan prácticas'', no ``personas''), \textbf{promesa} (qué resuelvo, en términos del beneficio que recibe el otro, no del esfuerzo que invierto yo) y \textbf{evidencia} (cómo se prueba, con ejemplos reales que ya existen, no con futuros condicionales).\\
2. Reconocimiento de patrones & \textbf{Reconocimiento de patrones.} Toda marca fuerte sigue el patrón ``Ayudo a [X] a [resolver/lograr Y] gracias a [Z]''. Una marca de panadería del barrio: ``Ayudo a familias de Cartago a celebrar lo importante con pan recién hecho gracias a la receta de mi abuela''. Cambia cualquiera de las tres variables y la marca pierde nitidez: si dices ``ayudo a personas'' pierdes audiencia, si dices ``con calidad'' pierdes el qué.\\
3. Abstracción & \textbf{Abstracción.} Lo esencial de tu marca cabe en \emph{una frase} --- si necesitas tres párrafos para explicar quién eres, todavía no la has abstraído. Esa frase es tu \emph{elevator pitch}: la que dirías en un ascensor, en una cita rápida, en la primera línea de tu LinkedIn. Si la tienes clara, todo lo demás (sitio, redes, perfil) se vuelve consecuencia natural. Si no, cualquier decisión visual o textual se sentirá improvisada.\\
4. Algoritmo conceptual & \textbf{Algoritmo.} Pasos para construirla: (1) lista 3 audiencias posibles --- lo más concretas que puedas (``estudiantes de comunicación que quieren agencia'', no ``jóvenes''); (2) lista 3 resultados que sabes producir, verificables con evidencia ya existente en tu vida; (3) combínalos en 3 frases con el patrón ``Ayudo a [X] a [Y] gracias a [Z]''; (4) léelas en voz alta --- la que no te avergüenza ni te infla es la auténtica. Esa es tu propuesta de valor.\\
\end{tabularx}
\normalsize

\begin{drawbox}{milcTurquesa}{Anatomía de una propuesta de valor sólida}
\textbf{Ayudo a} [audiencia concreta, no ``personas''] \textbf{a} [resolver/lograr resultado específico, no ``ayudar''] \textbf{gracias a} [tu fortaleza distintiva, no ``ser creativo'']. \\[2mm] \textbf{Ejemplo débil:} ``Soy creativo y me gusta el diseño.'' \\[1mm] \textbf{Ejemplo fuerte:} ``Ayudo a pequeños negocios de Cartago a contar su historia en video corto.''
\end{drawbox}

\begin{softbox}{milcMostaza}{milcGris}{Errores comunes y su corrección}
(1) Frases genéricas tipo ``soy proactivo y trabajador'': suenan a CV de los noventa, no diferencian a nadie. (2) Audiencia vaga (``personas'', ``jóvenes'', ``profesionales''): tan amplia que no le sirve a nadie en concreto; la audiencia útil cabe en una habitación. (3) Inflar logros que no sostienes con evidencia: el primer encuentro real expone la inflación. (4) Copiar la marca de alguien que admiras: tu autenticidad es lo único que no se puede copiar. (5) Esperar a tener ``cosas que mostrar'' para definir la marca: la marca se construye antes de la primera evidencia, no después --- es lo que ordena qué evidencia producir.
\end{softbox}

\begin{infoband}{milcTurquesa}


\textbf{Lo que va al cuaderno (Actividad 2 · EXPLICA).} Título: «Actividad 2 --- Mi propuesta de valor». \textbf{Formato:} 3 versiones de tu propuesta numeradas (1, 2, 3) siguiendo el patrón ``Ayudo a [X] a [Y] gracias a [Z]'', cambiando cada vez la audiencia o la fortaleza. Luego marca con $\star$ la versión elegida y escribe 2 renglones explicando por qué esa. \textbf{Sabes que terminaste cuando} cada frase nombra una audiencia concreta (no ``personas''), un resultado específico (no ``ayudar''), y la elegida la puedes decir en voz alta sin titubear.
\end{infoband}

\puente{Ya tienes tu propuesta de valor en una frase. Bien. Ahora la vas a convertir en un manifiesto completo --- una página de tu cuaderno que será tu ancla durante todo el periodo 1. Cada guía siguiente lo desarrolla; hoy lo dejas escrito.}

\begin{titlebox}{milcMagenta}
Fase 3 · Praxis con pensamiento computacional
\end{titlebox}

\begin{softbox}{milcMagenta}{milcGris}{Construyendo el producto con los cuatro pilares}
\textbf{Actividad 3 · CREA --- Tu Manifiesto de Marca Personal} (30 min · individual). En una página del cuaderno vas a escribir 5 bloques cortos con títulos fijos. Este manifiesto será tu ancla durante todo el periodo 1 --- cada guía siguiente lo desarrolla.
\end{softbox}

\small
\begin{tabularx}{\textwidth}{>{\bfseries\RaggedRight\arraybackslash}p{3.6cm}Y}
\rowcolor{milcMagenta!90}\color{white}Pilar & \color{white}\bfseries Aplicación al producto\\
Descomposición & \textbf{Descomposición.} Tu manifiesto se descompone en 5 bloques fijos, cada uno con función pedagógica distinta: \textbf{Quién soy} (te ancla en lo real, sin currículum), \textbf{A quién sirvo} (compromete a una audiencia concreta), \textbf{Qué prometo} (declara el valor entregado), \textbf{Qué NO soy} (delimita los rechazos --- crucial para tu identidad), \textbf{Cómo me reconocerán} (vuelve la promesa observable). Los 5 bloques actúan juntos; si falta cualquiera, la marca queda incompleta.\\
Patrones & \textbf{Patrones.} Sigue el patrón clásico de manifiestos fuertes: \textbf{lenguaje directo} (sin adjetivos vacíos como ``apasionado''), \textbf{primera persona} (yo, no ``el alumno''), \textbf{frases cortas} (no más de una idea por oración), \textbf{evidencia concreta} (un proyecto real, una pieza terminada). El opuesto se reconoce fácil: lleno de promesas vagas y sin un solo ejemplo verificable.\\
Abstracción & \textbf{Abstracción.} Cada bloque expresa una sola idea --- si un bloque tuyo dice tres cosas distintas, todavía no abstrajiste lo esencial. Es el momento más difícil del manifiesto: el primer borrador siempre dice de más. Releerlo al día siguiente con un lápiz para tachar es lo que separa una marca confusa de una nítida. El test final: léelo en voz alta a alguien que te conoce y mira si te reconoce.\\
Algoritmo de armado & \textbf{Algoritmo.} Pasos para no perderte: (1) escribe los 5 títulos en blanco en la página, dejando espacio entre cada uno; (2) llena ``Quién soy'' primero, en 2-3 renglones --- sin currículum, sin inflar; (3) cierra con ``Qué NO soy'' usando 3 cosas concretas (``no soy el que diseña carteles en una hora'', no ``no soy malo''); (4) ``Cómo me reconocerán'' con 3 evidencias observables (manera de escribir, temas que repites, una marca visual que ya tienes).\\
\end{tabularx}
\normalsize

\begin{drawbox}{milcMagenta}{Checklist del manifiesto antes de cerrar el cuaderno}
\checkbox Los 5 bloques están presentes y titulados. \\[1mm] \checkbox ``Qué NO soy'' tiene al menos 3 cosas concretas (no ``no soy malo''). \\[1mm] \checkbox ``Cómo me reconocerán'' da 3 evidencias observables (manera de escribir, temas que repites, una marca visual). \\[1mm] \checkbox Lo puedes leer en voz alta sin avergonzarte ni inflarte. \\[1mm] \checkbox Tu propuesta de valor (Actividad 2) se refleja de forma coherente.
\end{drawbox}

\begin{softbox}{milcMagenta}{milcGris}{Plantilla de trabajo}
\textbf{Quién soy.} \textit{En 2--3 renglones, quién eres hoy sin currículum.} \\[1mm] \textbf{A quién sirvo.} \textit{La audiencia concreta de tu marca en este momento.} \\[1mm] \textbf{Qué prometo.} \textit{El resultado que sabes producir.} \\[1mm] \textbf{Qué NO soy.} \textit{3 cosas con las que no quieres que te confundan.} \\[1mm] \textbf{Cómo me reconocerán.} \textit{3 evidencias observables --- no abstracciones.}
\end{softbox}

\begin{infoband}{milcMagenta}


\textbf{Lo que va al cuaderno (Actividad 3 · CREA).} Título: «Actividad 3 --- Manifiesto de marca personal · versión 1». \textbf{Formato:} 5 bloques titulados (Quién soy · A quién sirvo · Qué prometo · Qué NO soy · Cómo me reconocerán), 3--5 renglones cada uno. \textbf{Extensión:} una página completa de cuaderno. \textbf{Sabes que terminaste cuando} los 5 bloques están presentes con sus títulos exactos, ``Qué NO soy'' tiene al menos 3 cosas concretas, ``Cómo me reconocerán'' da 3 evidencias observables, y puedes leerlo en voz alta sin avergonzarte ni inflarte.
\end{infoband}

\puente{Lo que sigue es el resumen de \emph{qué} entregas hoy y \emph{cómo se sabe} si quedó bien hecho. Léelo antes de cerrar el cuaderno: son los criterios con los que el manifiesto se evalúa --- a ti mismo, primero.}

\begin{titlebox}{milcMostaza}
Producto de la sesión
\end{titlebox}
\textbf{Manifiesto de marca personal en una página de cuaderno}.
\begin{softbox}{milcMostaza}{milcGris}{Criterios de calidad}
(1) Los 5 bloques están presentes y titulados con el nombre exacto. (2) ``Qué NO soy'' tiene al menos 3 cosas concretas. (3) ``Cómo me reconocerán'' da 3 evidencias observables (no abstracciones). (4) El manifiesto suena a ti --- ni inflado ni copiado. (5) Tu propuesta de valor de la Actividad 2 se refleja con coherencia.
\end{softbox}

\puente{Antes de cerrar, mira tu trabajo desde cinco dimensiones humanas. Evaluar no es la nota: es lo que se queda en ti cuando esta guía pase. Personal, emocional, ciudadana, local, intergeneracional --- cinco preguntas pequeñas que cambian con el tiempo.}

\begin{titlebox}{milcVino}
Fase 4 · Evaluación liberadora — cinco dimensiones
\end{titlebox}

\begin{softbox}{milcVino}{milcGris}{Sentido de la evaluación}
Evaluar no es solo revisar una nota. Es reconocer qué aprendí, qué cambió en mí, qué emociones manejé, cómo me relaciono con la ciudad y la comunidad, y qué le dejaré a quien viene detrás.
\end{softbox}

\textbf{Mi reflexión liberadora en cinco dimensiones.} Completa cada frase iniciada:

\small
\begin{tabularx}{\textwidth}{>{\bfseries\RaggedRight\arraybackslash}p{3.4cm}Y>{\RaggedRight\arraybackslash}p{4.6cm}}
\rowcolor{milcVino}\color{white}Dimensión & \color{white}\bfseries Mi respuesta (frame starter) & \color{white}\bfseries Algo que descubrí\\
1. Personal & Lo que más cambió en mí fue \linead{6.5cm} & \linead{4.2cm}\\
2. Emocional & La emoción más fuerte fue \linead{2.5cm} y la regulé \linead{3cm} & \linead{4.2cm}\\
3. Ciudadana & En democracia esto importa porque \linead{6cm} & \linead{4.2cm}\\
4. Local (Cartago) & En mi barrio sería útil para \linead{6.5cm} & \linead{4.2cm}\\
5. Intergeneracional & A mi abuela/abuelo le diría \linead{6.5cm} & \linead{4.2cm}\\
\end{tabularx}
\normalsize

\begin{infoband}{milcVino}
\textbf{Para la próxima sustentación.} Vuelve a esta tabla en seis meses. Verás que las respuestas cambian — no porque lo aprendido cambie, sino porque tú cambias. La evaluación liberadora es una conversación contigo a través del tiempo.
\end{infoband}


\puente{Y para terminar, tres voces antiguas que te van a acompañar durante toda la secundaria. Dussel, un estoico, Floridi. No los memorices --- escúchalos. Cada uno te ofrece un lente distinto para mirar la marca personal que acabas de escribir.}

\begin{titlebox}{milcGradeDark}
Triángulo de pensamiento · cierre filosófico
\end{titlebox}

\begin{softbox}{milcVino}{milcGris}{Dussel · lente del nosotros}
\textit{``Sin reconocimiento del Otro no hay justicia, sólo poder.''}

Una marca personal puede ser un acto de reconocimiento o un acto de poder. Si tu marca solo te promociona a ti, ejerce poder; si abre espacio a otros, encuentra justicia. En el Valle, las marcas familiares más respetadas eran las que sostuvieron una comunidad.

\textbf{¿A quién sirve mi marca? ¿Reconozco al Otro o solo me promociono?}
\end{softbox}
\linead{\linewidth}\\[1mm]
\linead{\linewidth}

\begin{softbox}{milcMostaza}{milcGris}{Estoicismo · Marco Aurelio · cuidado interior}
\textit{``Tu vida es lo que tus pensamientos hagan de ella.''}

Lo que publicas, cómo respondes a un mensaje incómodo, cómo manejas un comentario hostil: todo \emph{eso} depende de ti. Los algoritmos, los juicios ajenos, el clima de las redes: \emph{eso no}. Distinguir las dos columnas es la ataraxia digital.

\textbf{¿La huella que dejo coincide con quien quiero ser? ¿Distingo lo que depende de mí de lo que no?}
\end{softbox}
\linead{\linewidth}\\[1mm]
\linead{\linewidth}

\begin{softbox}{milcTurquesa}{milcGris}{Floridi · lente de la infoesfera}
\textit{``Somos inforgs en una infoesfera que cuidamos o degradamos con cada acto informacional.''}

Cada publicación tuya es un acto informacional: añade valor a la infoesfera o la ensucia. Tu marca personal, en términos de Floridi, es una responsabilidad sobre el bien común informacional. No es vanidad: es ética digital aplicada.

\textbf{¿Mi marca enriquece o ensucia la infoesfera? ¿Lo que añado al mundo digital agrega valor a otros?}
\end{softbox}
\linead{\linewidth}\\[1mm]
\linead{\linewidth}

\begin{titlebox}{milcVino}
Compromiso final
\end{titlebox}

\small
\begin{tabularx}{\textwidth}{>{\RaggedRight\arraybackslash}p{3.0cm}YYY}
\rowcolor{milcVino}\color{white}\bfseries Criterio & \color{white}\bfseries Lo logré & \color{white}\bfseries En proceso & \color{white}\bfseries Necesito apoyo\\
Comprensión del tema & Explico ideas centrales con mis palabras. & Reconozco ideas pero falta precisión. & Necesito repasar conceptos base.\\
Pensamiento computacional & Apliqué los 4 pilares al diseñar mi producto. & Apliqué algunos pilares. & Necesito releer la fase 2.\\
Aplicación y producto & Mi producto cumple los criterios y se puede revisar. & Producto funcional con ajustes pendientes. & Aún en construcción.\\
Reflexión 5 dimensiones & Respondí cada dimensión con honestidad. & Respondí algunas con detalle. & Mis respuestas son muy generales.\\
Triángulo de pensamiento & Conecté las 3 voces con mi trabajo real. & Conecté al menos una. & No relaciono las voces con mi proceso.\\
\end{tabularx}
\normalsize

\textbf{Hoy entendí que:}\\[1mm]
\linead{\linewidth}\\[1mm]
\linead{\linewidth}

\textbf{Mi compromiso para la próxima sesión es:}\\[1mm]
\linead{\linewidth}\\[1mm]
\linead{\linewidth}

\begin{softbox}{milcMostaza}{milcGris}{Cita de despedida · Marco Aurelio}
\textit{``El obstáculo en el camino se vuelve el camino. Lo que estorba al camino se vuelve camino.''}\\[1mm]
Cada error de hoy, bien observado, se convierte en pieza del camino que pisarás mañana.
\end{softbox}

\small
\begin{tabularx}{\textwidth}{>{\bfseries}p{3.6cm}Y}
\rowcolor{milcGradeDark!12}Nombre completo: & \linead{10cm}\\
Fecha: & \linead{4cm} \quad Curso/grupo: \linead{4cm}\\
Mi firma: & \linead{9cm}\\
\end{tabularx}
\normalsize

\begin{infoband}{milcGradeDark}
\textbf{Para la próxima semana.} Antes de cada publicación en cualquier red, pregúntate: ¿esto suma o resta a la marca que escribí hoy en mi manifiesto? Si no suma, no se publica. Esa disciplina de una semana vale más que mil consejos. El próximo lunes traes tu manifiesto al aula y le añades los aprendizajes de la semana.
\end{infoband}

\end{document}
